从整体结果看，当前近似匹配方案表现出较强的攻击类型针对性，可归纳为以下规律。需要强调的是，该匹配表更适合用于提炼防守机制并支撑问题三的决策建模，不宜直接视作高保真仿真意义下的精确数值结论。

\analysisparagraph{1. 直线型拳法攻击优先采用快速格挡。}
对于刺拳和后手直拳，模型输出的前三防守分别为“单手拍挡(0.425)”“十字格挡(0.367)”“护头防御(0.345)”与“单手拍挡(0.405)”“十字格挡(0.352)”“左右侧闪(0.329)”。这说明面对轨迹清晰、正向突进的直线型拳法，低代价、短启动时间的拍挡动作最具性价比，而十字格挡与护头防御/侧闪可作为稳健型次优方案。

\analysisparagraph{2. 弧线型拳法更依赖专项防守。}
对于摆拳，模型给出的前三防守为“下潜闪避(0.366)”“肘挡(0.347)”“单手拍挡(0.340)”。这表明对大弧线、横向包络明显的攻击，单纯正面格挡并非最优，改变受击平面或采用更适合近身弧线攻击的肘挡更有效。

\analysisparagraph{3. 腿法攻击会诱导模型在“格挡”和“位移规避”之间做折中。}
对前蹬腿，前三防守为“单手拍挡(0.340)”“后撤步(0.317)”“下压格挡(0.298)”；对回旋踢，前三防守为“转身闪避(0.361)”“下潜闪避(0.341)”“单手拍挡(0.340)”。这说明当攻击范围扩大、动作恢复帧较长时，位移型闪避由于兼具防御效果与反击窗口，通常优于原地硬挡。

\analysisparagraph{4. 低段平衡破坏类攻击强调专项低位拦截。}
对于低扫腿，模型给出的前三防守为“下压格挡(0.398)”“单手拍挡(0.340)”“沉身防御(0.318)”。说明针对下肢攻击，模型不再只依赖一般性格挡能力，而是明显偏向更具针对性的低位防守动作。

\analysisparagraph{5. 倒地状态下的防守逻辑发生根本改变。}
在状态 $4$（Down）下，跨攻击统计中第一、第二推荐动作分别集中于“倒地防御”与“快速起身”。这说明一旦机器人失去站立能力，模型的决策目标将由主动拦截来袭攻击，切换为先降低地面追加伤害，再尽快恢复站立姿态。
